\documentclass[12pt]{article}
\usepackage[utf8]{inputenc}
\usepackage[T1]{fontenc}
\usepackage{hyperref}
\title{Marangis Language Resources for K-12 Curriculum Development}
\author{Compiled by Claude Sonnet 5 (claude-sonnet-5) for Andr\'as Kornai}
\date{}
\begin{document}
\maketitle
\section{Introduction}
Marangis, also called Watam (ISO 639-3: wax; Glottocode wata1253), is a Ramu language of Papua New Guinea, spoken in Marangis village (Madang Province) and Watam village (East Sepik Province). Ethnologue gave 590 speakers (2003); a later SIL survey reported about 1,086 community members and classed the language as threatened. Both communities said they want children to learn it.
\section{Primary Educational Resources}
None found.
\section{Linguistic and Reference Materials}
\begin{itemize}
\item Wikipedia: \url{https://en.wikipedia.org/wiki/Marangis_language}
\item Endangered Languages Project entry: \url{https://endangeredlanguages.com/lang/10645}
\item Endangered Languages Project source listing: \url{https://endangeredlanguages.com/node/98322}
\end{itemize}
\section{Religious and Community Resources}
None found.
\section{Digital and Audio Resources}
None found.
\section{Curriculum Development Context}
Little documentation was found. Community attitudes are favourable, but materials would have to be created.
\end{document}